\chapter{记号与规划模板}
\label{app:notation}

\section{记号表}

\begin{center}
\begin{tabularx}{\textwidth}{@{}l X@{}}
\toprule
符号 & 含义 \\
\midrule
$s$ & 行为状态 \\
$\FI(s)$ & 状态 $s$ 的框架惯性 \\
$\Reward(s)$ & 当前状态的即时奖励率 \\
$\Switch(s)$ & 离开当前状态的认知切换成本 \\
$\Env(s)$ & 当前状态的环境强化 \\
$\DV(s_0 \to s_1)$ & 一次转变的激活障碍 \\
$\Will(t)$ & 时刻 $t$ 可用的意志能量 \\
$\Leverage(t)$ & 时刻 $t$ 自由意志的有效杠杆 \\
$P$ & 用于降低未来障碍项的前置准备 \\
\bottomrule
\end{tabularx}
\end{center}

\section{规划模板}

\begin{templatebox}
\textbf{单次会话设计卡}
\begin{enumerate}[nosep]
  \item 本周的目标结果：
  \item 服务于该结果的下一次会话：
  \item 精确的第一个动作：
  \item 必须提前摆放好的物件：
  \item 这次会话应当开始的窗口：
  \item 结束这次会话的规则：
\end{enumerate}
\end{templatebox}

\begin{templatebox}
\textbf{失败复盘卡}
\begin{enumerate}[nosep]
  \item 哪个状态赢了？
  \item 哪个障碍项被低估了？
  \item 是哪个线索恢复了旧状态？
  \item 哪个更小的进入动作应当替换当前动作？
  \item 在下一次窗口到来之前，会先改变什么？
\end{enumerate}
\end{templatebox}

\section{附录用途}

本附录是为复用而存在的。如果后续版本加入更多形式化证明、
经验性引文或特定领域的案例库，这些模块仍然可以共享同一套记号和
规划模板。这样，手稿在扩展时仍能保持整体一致性。